% Level 2 of data access. The lookup chain (access.chain_for) decides where each
% file is read, one locator per place: dataset roots, staging, bundles, the
% restricted cache, a link or a copy in the public cache, a download. Bundles
% sit under it because the chain reads them, not beside Selection.
\begin{tikzpicture}
\node[actor, text width=34mm, minimum height=13mm] (selection) at (0,0) {Selection};
\node[actor, text width=34mm, minimum height=13mm] (retrieval) at (5,0) {Retrieval};
\node[actor, text width=34mm, minimum height=13mm] (access) at (10,0) {Lookup chain\\\texttt{access}};
\node[artifact, text width=34mm, minimum height=13mm] (staging) at (0,-2.7) {Staging};
\node[artifact, text width=34mm, minimum height=13mm] (integrity) at (5,-2.7) {Local integrity tools\\\texttt{verify / materialize}};
\node[artifact, text width=34mm, minimum height=13mm] (bundles) at (10,-2.7) {Repository bundles};
\node[artifact, text width=34mm, minimum height=13mm] (catalog) at (0,-5.5) {Catalogue model\\shared block};
\node[artifact, text width=34mm, minimum height=13mm] (pooch) at (5,-5.5) {Downloader port\\Pooch / filesystem};
\node[artifact, text width=34mm, minimum height=13mm] (config) at (10,-5.5) {Configuration\\shared block};
\draw[flow] (selection) -- node[edgelabel,right]{development overlay} (staging);
\draw[flow] (retrieval) -- node[edgelabel,above]{locations} (access);
\draw[flow] (selection.west) -- ++(-7mm,0) |- (catalog.west);
\draw[flow] (staging) -- node[edgelabel,right]{overlay} (catalog);
\draw[flow] (access) -- node[edgelabel,right]{bundled files first} (bundles);
\draw[flow] (bundles.south west) -- node[edgelabel,above,sloped]{export} (pooch.north east);
\draw[flow] (integrity.north east) -- node[edgelabel,above,sloped]{locations} (access.south west);
\draw[flow] (access.east) -- ++(7mm,0) |- (config.east);
\draw[flow] (retrieval.south west) -- (2.5,-1.2) -- (2.5,-4.3) -- (pooch.north west);
\node[note, rotate=90] at (2.25,-2.75) {verified fetch};
\begin{scope}[on background layer]
\node[boundarybox, fit=(selection)(access)(staging)(bundles), label={[boundarylabel]above:Data access -- level 2}] {};
\end{scope}
\node[note, text width=140mm] at (5,-7) {Consumer entry points coordinate Selection then Retrieval; the chain answers found, pass or refuse for each file.\\Bundle export selects canonical resources through Selection; supplementary dependencies are omitted.};
\end{tikzpicture}
